\chapter{Dataset}
\label{ch:dataset}

\section{Source and composition}

The corpus consists of short clips of Bangla speech cut from television shows,
films and online video, originally distributed as MP4 and grouped into one
folder per class. It contains nine classes: eight regional dialects named after
the district or division they are associated with, and the standard register,
labelled Formal (\cref{tab:classes}). After validation, 9,303 clips totalling
about 13.1 hours are usable; two duplicate \path{.ogg} files in the Formal
folder are excluded. Class sizes range from 750 (Khulna) to 1,652 (Formal).

\begin{table}[ht]
  \centering
  \caption{Classes and clip counts, with the number of test clips under the two
  split protocols of \cref{sec:method-splits}.}
  \label{tab:classes}
  \begin{tabular}{lrrr}
    \toprule
    Class & Clips & Test (clip-level) & Test (strict) \\
    \midrule
    Barishal   & 1,257 & 179 & 251 \\
    Chattogram &   797 & 114 & 114 \\
    Dhaka      &   763 & 109 & 103 \\
    Formal     & 1,652 & 236 & 236 \\
    Khulna     &   750 & 107 &  67 \\
    Mymensingh & 1,053 & 151 & 150 \\
    Noakhali   & 1,213 & 174 & 173 \\
    Rajshahi   &   860 & 122 & 110 \\
    Sylhet     &   958 & 137 & 123 \\
    \midrule
    Total      & 9,303 & 1,329 & 1,327 \\
    \bottomrule
  \end{tabular}
\end{table}

\section{Audio properties}

Every clip is standardised to 16\,kHz, mono, 16-bit PCM WAV
(\path{build_manifest.py}). Clips are nominally five seconds long: the median
duration is 5.0\,s and the mean 5.08\,s, but 1,067 clips fall outside
4.5--5.5\,s, with a range of 1.1--8.0\,s; Formal clips are slightly longer
(median 5.6\,s). Five seconds is enough for a listener to recognise many
dialects but excludes longer-range cues such as discourse markers and phrase-level
intonation. Because the material is broadcast and online media, recording
conditions vary widely within and across classes, from studio news to
outdoor drama scenes.

\section{The Formal class}

Formal is a speech register---the standard, prescribed Bangla of news reading
and formal speech---rather than a region. It is kept as a ninth class
deliberately. Distinguishing regionally marked speech from the standard register
is of practical interest in its own right, and Formal is the only class with
recoverable provenance, which makes it the calibration anchor for the strict
split (\cref{sec:method-splits}). A consequence is that the class boundary is
partly one of register: Formal speakers may come from any region, and some
regional clips contain stretches of near-standard speech.

\section{Metadata and the leakage problem}
\label{sec:data-metadata}

The corpus ships no speaker labels and almost no provenance. Regional clips have
sequential filenames (\path{khl271.wav}, \path{shlt479.wav}) that say
nothing about speaker or source. Only the Formal class carries recoverable
source identifiers: 557 of its 1,652 clips embed a YouTube video identifier in
the filename, from 35 distinct videos. For the remaining 8,746 clips---94\,\% of
the corpus---nothing prevents clips of the same speaker, or from the same
programme and recording channel, from appearing on both sides of a random
train/test split. \Cref{sec:method-splits} describes how the thesis estimates
pseudo-speaker groups to control this, and \cref{sec:res-strict} measures its
effect.

The clips were collected from films and dramas as well as from news and
interviews, so some speakers may be actors performing a dialect rather than
native speakers of it. This cannot be checked from the metadata and is listed
among the limitations (\cref{sec:disc-limits}).

\section{External resources for the lexical route}
\label{sec:data-external}

The lexical route (\cref{sec:method-lexical}) uses three public text and speech
resources in addition to the corpus. None contributes training audio to any
dialect classifier.

\begin{description}
  \item[OOD-Speech] \citep{rakib2023oodspeech}: the 963,636 training sentences
    of standard Bangla, distributed with the 2023 Bengali.AI competition
    \citep{bengaliai2023competition}, define the standard-Bangla vocabulary.
  \item[Ben-10] \citep{dipto2025ben10}: the district-labelled dialect
    transcripts of its training set (13,342 clips) contribute regional words for
    the three districts it shares with this corpus (Barishal, Chattogram,
    Sylhet); its public validation set (1,666 clips with audio and human
    transcripts) is the reference for measuring speech-recognition quality.
    Released under CC0.
  \item[Vashantor] \citep{faria2023vashantor}: regional--standard sentence
    pairs for Chittagong, Noakhali, Sylhet, Barishal and Mymensingh contribute
    regional words for those five classes. Released under CC BY 4.0.
\end{description}
